\documentclass[aspectratio=169,11pt]{beamer}
\usefonttheme{professionalfonts}
\renewcommand{\normalsize}{\fontsize{12}{15}\selectfont}
\usepackage[T1]{fontenc}
\usepackage{lmodern,amsmath,booktabs,tabularx,array,tikz}
\usetikzlibrary{arrows.meta,positioning}
\definecolor{navy}{HTML}{1D3550}
\definecolor{teal}{HTML}{147D83}
\definecolor{grey}{HTML}{596572}
\definecolor{rulegrey}{HTML}{D8DFE5}
\setbeamercolor{normal text}{fg=navy,bg=white}
\setbeamercolor{structure}{fg=teal}
\setbeamercolor{frametitle}{fg=navy,bg=white}
\setbeamerfont{frametitle}{size=\large,series=\bfseries}
\setbeamertemplate{navigation symbols}{}
\setbeamertemplate{itemize items}[circle]
\setbeamersize{text margin left=9mm,text margin right=9mm}
\setbeamertemplate{frametitle}{%
 \vspace{3mm}\insertframetitle\par\vspace{2mm}%
 {\color{rulegrey}\hrule height .4pt}\vspace{1mm}}
\setbeamertemplate{footline}{%
 \hspace{9mm}\begin{beamercolorbox}[wd=.86\paperwidth,ht=3ex,dp=1.5ex]{normal text}%
 \color{grey}\scriptsize Saarthak Krishan\hfill DH 307\hspace{5mm}\insertframenumber
 \end{beamercolorbox}}
\addtobeamertemplate{itemize/enumerate body begin}{\setlength{\itemsep}{4mm}}{}
\newcommand{\process}[3]{%
 \begin{center}
 \begin{tikzpicture}[node distance=10mm,every node/.style={font=\small\bfseries},every path/.style={draw=teal,line width=.8pt}]
 \node (a) {#1}; \node[right=of a] (b) {#2}; \node[right=of b] (c) {#3};
 \draw[-{Latex}] (a)--(b); \draw[-{Latex}] (b)--(c);
 \end{tikzpicture}
 \end{center}\vspace{2mm}}
\newcommand{\conclusion}[1]{\par\vspace{4mm}\noindent{\color{teal}\bfseries #1\par}}
\title{Protein--Ligand Design}
\subtitle{Cooperative and Conformation-Aware Search}
\author{Saarthak Krishan}
\institute{Supervised by Prof. Debabrata Maiti}
\date{}

\begin{document}
\begin{frame}[plain]
\vspace{6mm}{\small\color{grey}DH 307 Research Progress\par}
\vspace{7mm}{\bfseries\fontsize{25}{30}\selectfont Protein--Ligand Design\par}
\vspace{3mm}{\large Cooperative and Conformation-Aware Search\par}
\vspace{5mm}{\color{teal}\rule{24mm}{1pt}\par}
\vspace{5mm}Saarthak Krishan\\[1mm]
{\small Roll No. 22B3959}\\[4mm]
{\small Supervised by Prof. Debabrata Maiti}
\end{frame}
\begin{frame}{Generation and decision-making are separate}

\begin{columns}[T]
\column{.47\textwidth}\textbf{Molecular generators}\par\vspace{3mm}
\[q_\theta(L\mid P),\qquad q_\phi(P\mid L,S)\]
FLOWR proposes ligands. PocketGen supplies pocket-edit proposals. Their weights stay fixed.
\column{.47\textwidth}\textbf{Selection policies}\par\vspace{3mm}
\[\pi_i(a\mid o,\mathcal A)\]
PPO trains policies to choose among available candidates using molecular features and rewards.
\end{columns}
\conclusion{Our learned component selects structures supplied by the generators.}
{\small\color{grey}$L$: ligand; $P$: pocket; $S$: scaffold; $o$: observation; $\mathcal A$: candidate set.}

\end{frame}
\begin{frame}{One decision step: propose, select, evaluate, learn}

\process{Fresh proposals}{Two selections}{Pair evaluation}
\begin{enumerate}
\item The wrapper requests candidates from the pretrained models.
\item The ligand actor selects a molecule; the second actor selects its partner.
\item The evaluator returns feasibility and utility for that pair.
\item Save observations, actions, rewards and old policy predictions for PPO.
\end{enumerate}
\conclusion{The actors use cheap features before the selected pair is scored.}
{\small\color{grey}Partner = pocket edit in cooperation; protein conformation in competition.}

\end{frame}
\begin{frame}{The policy observes a 40-dimensional summary}

\begin{tabularx}{\textwidth}{@{}Xr@{}}
\toprule
Observation component & Channels\\\midrule
Ligand descriptors: size, weight, logP, TPSA, rotatable bonds, donors, acceptors, rings, QED & 9\\[2mm]
Protein amino-acid composition & 20\\[2mm]
Nearest-protein-atom contact-distance histogram & 8\\[2mm]
Current utility, remaining decision budget, ligand charge & 3\\\bottomrule
\end{tabularx}\par
\vspace{4mm}
\[o_t=\operatorname{encode}(P_t,L_t,U_t,b_t)\in\mathbb R^{40}\]
\conclusion{A compact observation can lose important geometric information.}
{\small\color{grey}It is not guaranteed to be a sufficient Markov state; raw atom coordinates are not policy inputs.}

\end{frame}
\begin{frame}{Actions are candidates, rather than atomic edits}

For each candidate $a_j$, construct a 40-channel feature vector $f_j$.
\[z_j=\operatorname{Actor}_{\theta_i}([o_t;f_j])\]
\[\pi_i(a_j\mid o_t,\mathcal A_t)=
\frac{m_j\exp(z_j)}{\sum_k m_k\exp(z_k)}\]
\begin{itemize}
\item $z_j$: candidate score; $m_j\in\{0,1\}$: eligibility mask.
\item Training samples an eligible index from this distribution.
\item Current ligand/pocket options can retain the existing design.
\end{itemize}
\conclusion{The same actor scores every candidate, allowing variable action counts.}
{\small\color{grey}If all actions are masked, the distribution is undefined; this needs explicit handling.}

\end{frame}
\begin{frame}{Each agent has its own actor and critic}

\begin{columns}[T]
\column{.47\textwidth}\textbf{Actor: candidate selection}\par\vspace{4mm}
\centering
State (40) + candidate (40)\par\vspace{2mm}
$\downarrow$\par
Linear $80\to128$ + Tanh\par
$\downarrow$\par
Linear $128\to64$ + Tanh\par
$\downarrow$\par
Linear $64\to1$\par\vspace{2mm}
Candidate logit
\column{.47\textwidth}\textbf{Critic: expected return}\par\vspace{4mm}
\centering
State (40)\par\vspace{2mm}
$\downarrow$\par
Linear $40\to128$ + Tanh\par
$\downarrow$\par
Linear $128\to64$ + Tanh\par
$\downarrow$\par
Linear $64\to1$\par\vspace{2mm}
$V_{\psi_i}(o_t)$
\end{columns}
\conclusion{Two actors and two critics; no shared neural weights between agents.}
{\small\color{grey}One Adam optimizer updates both agents. $V_{\psi}$ here denotes value, not the Vina score.}

\end{frame}
\begin{frame}{Cooperation: two choices, one improvement reward}

\[\pi(a_L,a_P\mid o)=\pi_L(a_L\mid o)\,\pi_P(a_P\mid o,a_L)\]
\begin{itemize}
\item Ligand actor selects a FLOWR candidate.
\item Pocket actor observes features using that ligand and selects a pocket option.
\item Feasible pairs update the current design; worsening moves are allowed in training.
\end{itemize}
\[r_L=r_P=U_{t+1}-U_t-0.02\]
\conclusion{The shared reward encourages coordinated improvement.}
{\small\color{grey}Infeasible pair: both rewards $=-2$. Backbone fixed; cumulative mutation cap = 4.}

\end{frame}
\begin{frame}{Competition: a ligand-aware conformation opponent}

\process{Select ligand}{Select difficult state}{Opposite rewards}
\begin{itemize}
\item Fixed-sequence conformations come from restrained molecular dynamics.
\item The opponent observes the chosen ligand and selects a training state.
\item Feasible ligand choices carry forward; the protein sequence is unchanged.
\end{itemize}
\[r_L=U(P_C,L),\qquad r_C=-U(P_C,L)\]
\conclusion{The aim is robustness over the allowed local conformations.}
{\small\color{grey}Infeasible pair: $r_L=-20$, $r_C=+20$. This can give the opponent an easy reward.}

\end{frame}
\begin{frame}{The evaluator supplies utility and feasibility}

\[U(P,L)=-V_{\mathrm{pose}}(P,L)-0.01E(L)-0.05N(P)\]
\begin{tabularx}{\textwidth}{@{}p{24mm}>{\raggedright\arraybackslash}X@{}}
$V_{\mathrm{pose}}$ & Vina score at the supplied pose; more negative is better.\\[3mm]
$E(L)$ & MMFF94s ligand energy before minus after local minimization.\\[3mm]
$N(P)$ & Residue identities changed relative to the starting protein.\\
\end{tabularx}\par\vspace{3mm}
Chemical, geometric and mutation constraints determine feasibility.
\conclusion{Utility measures pair quality; reward defines the learning objective.}
{\small\color{grey}The weights are heuristic. Utility is not experimentally measured binding affinity.}

\end{frame}
\begin{frame}{An episode is a short design trajectory}

\process{Reset to reference}{Up to four steps}{Store the trajectory}
\begin{itemize}
\item Choose a training system; initialize its reference pair and utility.
\item Each step generates proposals, samples actions and evaluates one pair.
\item The current design changes only when the selected pair is feasible.
\item Reset molecular state at the next episode; retain learned network weights.
\end{itemize}
\conclusion{An episode is one design attempt, not one complete dataset pass.}
{\small\color{grey}A proposal failure can terminate early. Flow-integration steps are separate from RL steps.}

\end{frame}
\begin{frame}{Returns and advantages: what was better than expected?}

For each player, accumulate rewards backward through its episode:
\[G_t^{(i)}=r_t^{(i)}+\gamma G_{t+1}^{(i)},\qquad \gamma=1.\]
Compare the observed return with the old critic prediction:
\[A_t^{(i)}=G_t^{(i)}-V_{\psi_i,\mathrm{old}}(o_t).\]
\begin{itemize}
\item Positive advantage: outcome exceeded the critic's expectation.
\item Negative advantage: outcome fell below that expectation.
\item Implementation normalizes advantages across the combined update buffer.
\end{itemize}
\conclusion{The code uses episode returns, with no GAE or critic bootstrapping.}

\end{frame}
\begin{frame}{Why a policy gradient changes selection probabilities}

The objective is expected return under the current selection policy:
\[J(\theta)=\mathbb E_{\tau\sim\pi_\theta}\!\left[\sum_t r_t\right].\]
A policy-gradient estimator uses
\[\nabla_\theta J\;\approx\;
\frac{1}{M}\sum_t A_t\nabla_\theta\log\pi_\theta(a_t\mid o_t,\mathcal A_t).\]
\begin{itemize}
\item Increase probability for actions with positive advantage.
\item Decrease probability for actions with negative advantage.
\end{itemize}
\conclusion{No gradient through Vina or either molecular generator is required.}
{\small\color{grey}For each player's update, the other player's behavior is part of the sampled environment.}

\end{frame}
\begin{frame}{PPO limits the incentive for large policy changes}

Compare the new and rollout probabilities of the sampled action:
\[\rho_t(\theta)=\frac{\pi_\theta(a_t\mid o_t,\mathcal A_t)}
{\pi_{\mathrm{old}}(a_t\mid o_t,\mathcal A_t)}.\]
\[L_{\mathrm{clip}}=\mathbb E\!\left[\min\left(
\rho_t A_t,\;\operatorname{clip}(\rho_t,1-\epsilon,1+\epsilon)A_t
\right)\right],\quad \epsilon=0.2.\]
\begin{itemize}
\item Stored old log-probabilities remain fixed during the update.
\item Recompute new probabilities on the same saved candidates and masks.
\item Clipping removes extra reward for overly aggressive changes in the favored direction.
\end{itemize}
\conclusion{Clipping shapes the objective; it is not a hard probability bound.}

\end{frame}
\begin{frame}{A numerical example of PPO clipping}
\small
Suppose the old probability of a selected action was $0.20$.\par
\vspace{2mm}
\begin{tabularx}{\textwidth}{@{}Xrr@{}}
\toprule
 & Positive advantage & Negative advantage\\
 & $A=+1$ & $A=-1$\\\midrule
New probability & $0.30$ & $0.10$\\[1mm]
Ratio $\rho$ & $1.50$ & $0.50$\\[1mm]
Unclipped term $\rho A$ & $1.50$ & $-0.50$\\[1mm]
Clipped term & $1.20$ & $-0.80$\\[1mm]
PPO uses the minimum & $1.20$ & $-0.80$\\\bottomrule
\end{tabularx}\par
\conclusion{Both examples limit further incentive for a large favored change.}
{\small\color{grey}Illustrative numbers; not measured action probabilities from the runs.}

\end{frame}
\begin{frame}{The full loss trains selection, value and exploration}

For each saved decision, the minimized loss is
\[\mathcal L_t=-L_{\mathrm{clip},t}
+0.5\big(V_{\psi_i}(o_t)-G_t\big)^2
-0.01\mathcal H\!\left(\pi_i(\cdot\mid o_t,\mathcal A_t)\right).\]
\begin{itemize}
\item \textbf{Policy loss:} improve the probability of useful decisions.
\item \textbf{Value loss:} predict the player's observed return.
\item \textbf{Entropy bonus:} discourage premature collapse onto one action.
\end{itemize}
\conclusion{Average losses across both agents; backpropagate through their MLPs.}
{\small\color{grey}Separate player returns and critics; shared optimizer and pooled advantage normalization.}

\end{frame}
\begin{frame}{Training schedule used in the implementation}

\begin{tabularx}{\textwidth}{@{}Xr@{}}
\toprule
Configuration & Value\\\midrule
Episodes per seed & 512\\[1mm]
Decision horizon per episode & Up to 4\\[1mm]
Collect experience before an update & 8 episodes\\[1mm]
PPO epochs per update & 4\\[1mm]
Adam learning rate & $3\times10^{-4}$\\[1mm]
Ratio clip / entropy coefficient & $0.2\;/\;0.01$\\[1mm]
Gradient-norm clip & $1$\\\bottomrule
\end{tabularx}\par
\conclusion{64 rollout-update cycles per seed; three seeds per setting.}
{\small\color{grey}Each PPO epoch uses the full buffered batch. Clear the buffer after the update.}

\end{frame}
\begin{frame}{Evaluation uses learned ranking and explicit screening}

\begin{columns}[T]
\column{.47\textwidth}\textbf{During training}\par\vspace{3mm}
\begin{itemize}
\item Sample actor actions.
\item Score the selected pair.
\item Save experience and update PPO.
\end{itemize}
\column{.47\textwidth}\textbf{During evaluation}\par\vspace{3mm}
\begin{itemize}
\item Fix the validation-selected policies.
\item Rank candidates for scoring.
\item Retain the best feasible design found.
\end{itemize}
\end{columns}
\vspace{4mm}
In competition, final selection uses exact worst-case screening over training states;
held-out conformations are assessed afterward.
\conclusion{The opponent is a training component, not the final evaluation oracle.}

\end{frame}
\begin{frame}{Completed cooperative results: a small learning effect}

\noindent{\small Three completed policy seeds; higher mean utility is better.\par}
\vspace{4mm}
\begin{tabularx}{\textwidth}{@{}Xr@{}}\toprule
Method & Mean utility\\\midrule
Learned cooperative policy & 6.350\\[2mm]
Random joint search & 6.301\\[2mm]
Greedy alternating & 6.227\\[2mm]
Ligand-only & 5.902\\\bottomrule
\end{tabularx}\par
\vspace{3mm}
The paired uncertainty interval for learned minus random includes zero.
Ligand-only is limited by candidate eligibility.
\conclusion{We have not established a reliable advantage from RL.}

\end{frame}
\begin{frame}{Bottlenecks that PPO alone cannot resolve}

\begin{itemize}
\item \textbf{Candidate support:} the policy cannot select useful structures that are never proposed or eligible.
\item \textbf{Representation:} coarse descriptors lose local contact chemistry.
\item \textbf{Imbalance:} effective pocket and ligand action spaces differ; their contributions need measurement.
\item \textbf{Competitive feedback:} infeasible pairs can reward the opponent without teaching useful robustness.
\end{itemize}
\conclusion{Completed training is not evidence of physical validity or game convergence.}

\end{frame}
\begin{frame}{Next experiments: distill, refine and isolate learning}

\process{Feasible search teacher}{Imitate its choices}{Refine with PPO}
\begin{enumerate}
\item Align generation and eligibility; explicitly handle empty action sets.
\item Use teacher decisions to initialize candidate-selection policies.
\item Test a balanced curriculum or extra updates for the weaker selector.
\item Compare imitation, PPO, and imitation + PPO against simple search.
\end{enumerate}
\conclusion{Teacher quality, proposal coverage and matched budgets remain essential.}
{\small\color{grey}These interventions are proposed experiments; the current selectors were trained with PPO.}

\end{frame}
\begin{frame}[plain]
\vfill{\bfseries\fontsize{28}{34}\selectfont Thank you\par}
\vspace{4mm}{\color{teal}\rule{24mm}{1pt}\par}
\vspace{4mm}Saarthak Krishan\vfill
\end{frame}
\appendix
\begin{frame}{Backup: sources and implementation mapping}

\small
\textbf{PPO:} Schulman et al., Proximal Policy Optimization Algorithms.\par
\url{https://arxiv.org/abs/1707.06347}\par\vspace{3mm}
\textbf{FLOWR:} ligand proposals.\par
\url{https://arxiv.org/abs/2504.10564}\par\vspace{3mm}
\textbf{PocketGen:} pocket proposals; substitute for the PocketFlow direction.\par
\url{https://doi.org/10.1038/s42256-024-00920-9}\par\vspace{3mm}
\textbf{Local implementation:}\par
\texttt{research/models.py}: actors, critics, PPO loss.\par
\texttt{research/features.py}: observations and masks.\par
\texttt{research/run.py}: episodes, rewards and evaluation.\par
\texttt{research/evaluator.py}: pair utility and feasibility.

\end{frame}
\end{document}
